% page 559 du fac-simile (image p-558 du scan)
% bloc 170.2 x 233.3 mm ; marge gauche 31.8 mm
\pgc{10.160mm}{0.000mm}{
\l{\cel{55}551}
\l{}
\l{\sou{stativo}. (\sou{fotogr}.)Sorto di suportilo kun tri pedi, ek tri parti}
\l{\cel{3}faldebla, uzata por la fotografili di la fotografisti. - D}
\l{}
\l{\sou{statolito}. (bot.) Nomo di la amilo-grani qui, per lia pezo, kon-}
\l{\cel{3}tributas igar la vejetanti sentiva pri la efiko da gravito.}
\l{\cel{3}(\sou{geotropismo}). -}
\l{}
\l{\sou{statoro}. (\sou{elektro}). (en\cel{1}\sou{motoro kun elektrofluo polifaza}). Parto}
\l{\cel{3}fixa di la mashino, opoze a\cel{1}\sou{rotoro}. -}
\l{}
\l{\sou{statuo}. Figuro skultura, generale ek marmoro, ma anke ek metalo,}
\l{\cel{3}ek ligno, qua reprezentas tota-reliefe homo od animalo. - DEFIRS}
\l{}
\l{\sou{staturo}. Longeso di la korpo di homo, mezurata kande lu stacas.}
\l{\cel{3}- DEFIS.}
\l{}
\l{\sou{statuto}. (\sou{yurocienco}).To quo rezolvesis kom aplikenda (pri la}
\l{\cel{3}personi, la kozi, en ta o ca kazo). DEFIRS.}
\l{}
\l{\sou{stearino}. (\sou{kemio}). Materio ek lameli, perlomatrea, ek glicerino}
\l{\cel{3}e CH3.(CH2)16 CO2 H, quan onu extraktas de sebo. - DEFIRS.}
\l{}
\l{\sou{steatito}. (\sou{mineral.}) Silikato di magnezio, variatato di talko}
\l{\cel{3}kompakta, granetatra. - DEFIS.}
\l{}
\l{stebar. (\sou{sutado}). Trairar steke stofo per du fili quin onu sinkas,}
\l{\cel{3}per la pinto di agulo suto-mashinala, aden trui di qui singla}
\l{\cel{3}ek le sequanta situesas egaldiste de la pre-esanta. -}
\l{}
\l{\sou{stekar}. (\sou{trans.)}\cel{2}Penetrigar per koakto la pinto di objekto (kom}
\l{\cel{3}ex.: paliso aden tero, e c.) DE.}
\l{}
\l{\sou{stelo}. I. (\sou{linguo ne-ciencala)}. Omna astro qua brilas en firmamen-}
\l{\cel{3}to (ecepte Suno e Luno). - II. (\sou{linguo ciencala}) Korpo cielala,}
\l{\cel{3}lumifanta, sen movo videbla, e di qua la lumo quaze cintilifas.}
\l{\cel{3}- III. (\sou{astrologio)}.\cel{2}Astro qua, acensante en cielo lor la nasko}
\l{\cel{3}di homo, egardesas kom influanta bone o male lua fato. - IV.}
\l{\cel{3}\sou{(teatro, cinemo).}\cel{2}Artisto qua brilas, unika o preske unika, en}
\l{\cel{3}trupo de aktori. - V. (\sou{metaf.}) (pro analogio kun la radiifo di}
\l{\cel{3}la\cel{1}\sou{steli}) (1) punto briloza ; (2) objekto di qua la parti dispo-}
\l{\cel{3}zesas quale la spoki di roto, e qua memorigas la formo di stelo}
\l{\cel{3}reprezentita desegne. - EFIS.}
\l{}
\l{stelario. (\sou{bot.)}\cel{2}Planto unyardura, ek la familio \textquotedbl{}primulacei\textquotedbl{},}
\l{\cel{3}quan konsumas ula uceli. - L. stellaria.}
\l{}
\l{steleo. (arkitekt.) (epoki antiqua).\cel{2}Monolito di qua la formo}
\l{\cel{4}esas olta di la fusto di kolono, di kolono ruptita o di cipo,}
\l{\cel{4}destinita (maxim-multa-kaze) recevor epigrafo. - DEFI.}
\l{}
\l{\cel{1}+stejo. La parto di la teatro ube pleas la aktori. (\sou{anke metaf.})}
\l{\cel{4}- E.}
}
